% Transcription — fonds Alexandre Grothendieck, shelfmark 64, batch 5 (pages 81-100).
% Established from the facsimile archives/batches/64/batch-05.pdf (a copy of 64.pdf
% fetched through the project relay, no cover sheet: PDF page k = archivists' page 20(N-1)+k),
% per the skill transcribe-grothendieck. The folder is 156 pages in eight batches.
%
% Pass: Opus 5.5 (claude-opus-5-5), 2026-09-27 — first pass, unchecked against the pages by a human.
%
% Pages skipped: none as unrelated. Page 100 is a blank sheet carrying only a heading in his
%   hand ("IV Descriptions explicites par fonctions P de Weierstrass"); per hand.md it gets no
%   \page{} and its words become the closing \section, which opens the run of the next batch.
% His pagination: a run numbered on every page, 5 to 13 on archive pp. 81, 83, 85, 87, 89,
%   91, 93 (p. 81 = his 5, ..., p. 93 = his 11); the even pages carry no number of his, and
%   pp. 94-99 none either (as far as the tiles show). Equations numbered (30) to (78), continuing
%   a numbering begun in an earlier batch ((5), (27) are cited).
% What the batch is about: Legendre "épinglages" (rigidifications) of a relative elliptic curve
%   E/S in char. != 2, via the cubic étale cover J = (2E)^* and the conic/P^1 Sigma_J with its
%   divisor D_J; the torsor of Legendre épinglages under mu_4 and its mu_2-quotient
%   mu_4^* ^ omega_J; Legendre-Jacobi level-2 épinglages and the fine moduli scheme
%   U_{0,3} x mu_4^*; the action of S_3 x mu_4 (the group D'_3 = S_3 x_{Z/2} Z/4) on it and on
%   O(1); then a start on level-4 Jacobi épinglages, and a fully explicit computation over an
%   algebraically closed field k (V_J = Ker(k^J -> k), F_t, the families psi_i^omega(delta)).
% Notation fixed once for the batch: \underline{J} for the cubic cover (his underlined J);
%   \underline{\omega}_{\underline{J}} for the two cyclic orders; \mu_4^{*} for his double-stemmed
%   mu_4 with a star (primitive 4th roots of unity); \Sigma_{\underline{J}}, \Sigma_{\underline{J}}^{*};
%   \underline{O}; \mathbb{D}_3, \mathbb{D}'_3 for his double-barred D; ^{\natural} for an
%   unidentified lightning-shaped superscript on O(1); \mathcal{L} for his script L; \mathfrak{S}.
% \ill{}: 42. \uncertain{}: 36.

\documentclass[11pt,a4paper]{article}
\input{../preamble/grothendieck}

\folder{64}
\batch{5}
\pages{81}{100}
\dating{[à partir de 1980- 1982]}
\watermark{Édition de démonstration}
\foldertitle{Modules courbes elliptiques en caractéristiques ≠2 : notes manuscrites (s.d.)}

\begin{document}

\page{81}\note{p.~5 de l'auteur, chiffrée en tête de page ; la page s'ouvre au
milieu d'un développement commencé dans le lot précédent (formules (27) et
suivantes).}
\[
  \text{(30)} \qquad T_{E/S,0}^{\otimes 2} \simeq T_{\Sigma_{\underline{J}},t}
\]
où $t$ est la section de $\Sigma_{\underline{J}}^* = \Sigma_{\underline{J}} \setminus D_{\underline{J}}$
\uncertain{image} de la section nulle de $E$. Ceci posé, on définit un
\underline{épinglage de Legendre} de $E/S$ comme étant \struck{formé}
\uncertain{la donnée} d'un iso
\[
  \text{(31)} \qquad \underbrace{\Delta_{E/S}}_{T_{E/S,0}}
  \xrightarrow[\;\sim\;]{k} \underline{O}(1)_t \simeq
\]
tel que \struck{il existe un} l'isom $k^{\otimes 2}$ déduit
par passage au carré, \uncertain{lorsqu'on} \uncertain{le fait suivre} de l'identification,
via (30) et (27), à l'isom
\[
  \text{(32)} \qquad T_{\Sigma_{\underline{J}},t} \xrightarrow{\;\sim\;}
  T_{\Sigma_{\underline{J}},t} \wedge^{\{\pm 1\}}
  \bigl(\mu_4^{*} \wedge \underline{\omega}_{\underline{J}}\bigr)
\]
déduit d'une section
\[
  \text{(33)} \qquad
  \begin{cases}
    \alpha \in \Gamma\bigl(S, \mu_4^{*} \wedge \underline{\omega}_{\underline{J}}\bigr)
      & \text{i.e.\ d'un iso} \\[2pt]
    \alpha : \mu_4^{*} \simeq \underline{\omega}_{\underline{J}}
  \end{cases}
\]
(qui est déterminé de façon unique en termes
\note{dans (31), un symbole biffé précède $\Delta_{E/S}$ et un autre le
second membre ; l'exposant de $\underline{O}(1)$ est surchargé
(\uncertain{un $\otimes$} biffé) et n'est pas lu, et la formule s'arrête
sur « $\simeq$ » sans second membre. Le symbole que nous rendons
$\mu_4^{*}$ est un $\mu$ à hampe double, d'indice $4$ (lu nettement
page~82) et d'exposant $*$ : l'ensemble des racines primitives quatrièmes
de l'unité, torseur sous $\{\pm 1\}$.}

\page{82}
de $k$). On peut dire aussi que la donnée de $k$ équivaut à celle d'un
\struck{paire} \add{couple} $(\alpha, k)$, où $\alpha$ est une section de
$\mu_4^{*} \wedge \underline{\omega}_{\underline{J}}$, et où $k$ est un iso
(31) \underline{compatible} avec $\alpha$ en un sens évident. Pour $\alpha$
fixé, les $k$ correspondants forment les sections d'un torseur sous
$\mu_2$ --- le \underline{torseur des épinglages de Legendre} de $E/S$,
\struck{\ill{} de} \add{\uncertain{associé à}} « l'épinglage partiel » $\alpha$.

La catégorie \struck{\uncertain{groupoïde}} \add{fibrée} (qui est un champ en
groupoïdes) des courbes elliptiques relatives, munies d'un épinglage de
Legendre, équivaut à celle des triplets $(\underline{J}, t, \alpha)$, où
\[
  \begin{cases}
    \text{a) } \underline{J} \text{ rev.\ étale cubique de } S \\
    \text{b) } t \text{ section de } \Sigma_{\underline{J}}^* \\
    \text{c) } \alpha \text{ section de } \mu_4^{*} \wedge \underline{\omega}_{\underline{J}}
      \text{, i.e.\ iso.\ } \mu_4^{*} \simeq \underline{\omega}_{\underline{J}}
  \end{cases}
\]
\marginal{\underline{NB} Il faudrait expliciter la construction de $E$, en
fonction des données \ill{} (\uncertain{via} \uncertain{fonction} $\wp$ de
Weierstrass) \ill{}}
\note{aucun numéro n'est écrit devant ces données ; la page~83 renvoie
à « un objet (34) », qui semble être un tel triplet. La note
marginale est écrite en oblique dans le coin inférieur gauche, et
se rattache aux données a), b), c) par une accolade.}

\page{83}\note{p.~6 de l'auteur.}
\struck{On définit des}
Le \struck{catégorie} \add{champ} en question n'est pas rigide --- i.e.\ un
objet peut avoir des automorphismes non triviaux.
\struck{Les} \add{\uncertain{Les}} automorphismes \add{d'un objet (34)} \struck{sont} \add{\ill{}} définis par des
automorphismes de $\underline{J}$, i.e.\ \struck{deux} des sections des
schémas en groupes étales $\mathfrak{S}_{\underline{J}}$, qui sont
compatibles avec les données b) et c). La compatibilité avec c)
\underline{signifie} que c'est une section de
\[
  \mathfrak{S}_{\underline{J}}^{+} \simeq
  \underbrace{S}_{\substack{\text{correspond}\\ \text{à la section}\\
  \text{unité des}\\ \text{schémas en groupes}}}
  \amalg\ \underline{\omega}_{\underline{J}}
\]
une section qui est $\neq \mathrm{id}$ sur tte fibre géom.\ revient donc
: une section de $\underline{\omega}_{\underline{J}}$, soit $\rho$ --- de
sorte que $\underline{J}$ devient un torseur sous
$(\mathbb{Z}/3\mathbb{Z})_S$. La compatibilité avec b)
\underline{signifie} que $t = \rho(t)$ i.e.\ que $t$ \struck{correspond à}
\add{est une} section du sous-schéma
\[
  \underline{P}_{\underline{J}} \subset \Sigma_{\underline{J}}^*,
  \qquad \underline{P}_{\underline{J}} =
  \bigl(\Sigma_{\underline{J}}\bigr)^{\mathfrak{S}_{\underline{J}}^{+}}
\]
\struck{\ill{} dès que \ill{}}
\note{un symbole biffé devant $S$ dans la première formule. Dans la
deuxième phrase, l'ajout interligne « \ldots\ d'un objet (34) \ldots » remplace
le début biffé ; sa lecture est incertaine. La
formule $\underline{P}_{\underline{J}} = (\Sigma_{\underline{J}})^{\mathfrak{S}_{\underline{J}}^{+}}$
est écrite sous la ligne et entourée.}

\page{84}
donc que $E$ est \struck{la} \add{une} courbe \add{elliptique} exceptionnelle
\uncertain{hexagonale} sur $S$\ldots

Pour éliminer les automorphismes, il est naturel de combiner l'épinglage de
Legendre avec un épinglage de \struck{\uncertain{Fonch}} \add{Jacobi} d'échelon~2, on
trouve l'épinglage de Legendre--Jacobi d'échelon~2, défini par un couple
$(k, u)$, où $k$ est un ép.\ de Legendre (31), et où $u$ est
\[
  \text{(35)} \qquad u : \underline{J} \xleftarrow{\;\sim\;} \{0, 1, \infty\}_S
\]
une trivialisation de $\underline{J}$, i.e.\ un isom
\[
  \text{(36)} \qquad \Sigma_{\underline{J}} \xleftrightarrow{\;\sim\;} \mathbb{P}^1_S
\]
transformant le diviseur $\{0\}_S + \{1\}_S + \{\infty\}_S$ en
$D_{\underline{J}}$. La catégorie \add{fibrée} des courbes elliptiques munies
d'un tel épinglage est alors rigide, et équivalente à celle définie par les
couples $(t, \alpha)$, où $t$ est une section de $U_{0,3}$, $\alpha$ une section
de $\mu_4^{*}$ --- donc la \struck{cat.} champ \uncertain{en}
\note{« hexagonale » est une lecture incertaine : le mot commence par un
\emph{h} et finit en \emph{-gonale} ; c'est ce que le contexte appelle (les
sections de $\underline{P}_{\underline{J}}$, points fixes de
$\mathfrak{S}_{\underline{J}}^{+}$, correspondent à la courbe d'invariant
$j = 0$). Dans (35), la flèche est tracée vers la gauche ; dans (36) elle a
une pointe à chaque bout. Le numéro de (36) surcharge un autre chiffre.}

\page{85}\note{p.~7 de l'auteur.}
groupoïdes envisagé est représentable par le schéma
\[
  \text{(37)} \qquad U_{03} \times \mu_4^{*}
\]
Le fibré $\Delta_E$ correspondant à la courbe elliptique épinglée
(\uncertain{par} Legendre--\ill{}) universelle sur la base
$U_{03} \times \mu_4^{*}$, est isomorphe canoniquement à
\struck{\ill{}}
\[
  \text{(38)} \qquad \Pi_0 = \underline{O}_{U_{0,3} \times \mu_4^{*}}(1)^{\natural}
\]
or ici, puisqu'on est \struck{sur} \add{au-dessus de} $\mu_4^{*}$, le fibré en
\uncertain{carrés} $Q_{0S}$ ($S = U_{03} \times \mu_4^{*}$) est trivial, d'où un iso.\
can.
\[
  \text{(39)} \qquad \Pi_0 \simeq \underline{O}_{U_{0,3}}(1)
\]
\struck{La question \ill{} certaine pour les définitions}
\add{Mais il faut expliciter le « \ill{} »} naturel pour ces définitions de
\struck{et les groupes} \ill{} épinglées, et examinons \add{son} action sur
$U_{03} \times \mu_4^{*}$ ainsi que sur $\Pi_0$.
\note{avant (38), une première écriture de la formule, où se lisent
$U_{0,3}$, $\Pi_0 =$, $\underline{O}_{U_{0,3}}(1)$ et $\times \mu$, est biffée.
L'exposant que nous rendons $^{\natural}$ (ici et pages~86, 93 et 96) est un
signe en zigzag, en forme d'éclair, que nous ne savons pas identifier. Le second nom
propre de la parenthèse n'est pas lu ; c'est le même tracé qui, biffé,
précède « Jacobi » page~84.}

\page{86}
Considérons une courbe elliptique relative $E/S$, alors les
\struck{données} \add{épinglages} de Legendre
$k : \Delta_{E/S} \simeq \underline{O}_{\Sigma}(1)_t^{\natural}$ forment les
sections d'un torseur \add{$\mathcal{L}_{E/S}$} sous $\mu_{4S}$ --- le
\underline{torseur des épinglages de Legendre} de $E/S$. Construction de
torseur par multiplication scalaire de $\mu_4 \subset (\mathbb{G}_m \ldots)$.
On \struck{peut le} \add{peut} décrire ce torseur aussi de la façon
suivante, sans faire intervenir explicitement
$\underline{O}_{\Sigma}(1)_t^{\natural}$, mais seulement
$\underline{O}_{\Sigma}(1)_t$. Considérons les isom
\[
  \text{(40)} \qquad k' : \Delta_{E/S} \simeq \underline{O}_{\Sigma}(1)_t
\]
tels que $k'^{\otimes 2}$, interprété comme un
\[
  \text{(41)} \qquad T_{\Sigma,t} \simeq T_{\Sigma,t} \wedge^{\{\pm 1\}}
  \{\underline{\omega}_{\underline{J}}\}
  \quad \text{(évidemment unique)}
\]
\add{un isom.} provienne d'une section de $\underline{\omega}_{\underline{J}}$.
\add{Les isom.} forment \struck{un torseur} \add{les sections} d'un torseur
sous $\mu_{4S}$, soit $Q_{E/S}$ (qu'on peut interpréter comme le schéma des
\uncertain{sommets} d'un
\note{dans le premier isomorphisme, un indice biffé après $\Sigma$. Dans
(41), un $k'^{\otimes 2}$ biffé précède $T_{\Sigma,t}$, dont l'indice porte
aussi une surcharge ; un symbole biffé devant $\underline{\omega}_{\underline{J}}$.
L'ordre des ajouts « un isom. » et « Les isom. » autour de (41) est
incertain.}

\page{87}\note{p.~8 de l'auteur.}
schéma en carrés combinatoires sur $S$) dont le schéma des diagonales
s'identifie à $\underline{\omega}_{\underline{J}}$ :
\[
  \text{(42)} \qquad \underline{Q}_{E/S}/\mu_2 \simeq \underline{\omega}_{\underline{J}}
\]
Ceci posé, on voit que l'on a un isom \add{can.} de $(\mu_4)_S$-torseurs
\[
  \text{(43)} \qquad Q_{E/S} \wedge^{\mu_4} Q_{0S} \simeq \mathcal{L}_{E/S}
\]
i.e.
\[
  \mathcal{L}_{E/S} \simeq \underline{\mathrm{Isom}}_{\mu_{4S}\text{-tors.}}
  \bigl(Q_{0S}^{-1}, Q_{E/S}\bigr)
\]
On aurait pu définir les épinglages de Legendre de $E/S$ comme les
trivialisations de $Q_{E/S}$, au lieu de celles de $\mathcal{L}_{E/S}$ ---
mais on aurait alors une relation moins \ill{} avec les épinglages de Jacobi
d'échelon~4.

On voit \add{donc} que les épinglages \add{de $E/S$} de Legendre forment
\add{les sections d'} un torseur sous $\mu_4$, telles que le $\mu_2$-torseur
associé soit can.\ iso.\ \add{à $\underline{\omega}_{\underline{J}} \wedge \mu_4^{*}$} ;
\[
  \text{(44)} \qquad \mathcal{L}_{E/S}/\mu_2 = \mathcal{L}_{E/S}
  \wedge^{\mu_{4S}} \mu_{2S} \simeq
  \bigl(\mu_{4S}^{*} \wedge \underline{\omega}_{\underline{J}}\bigr)
\]
tandis que les épinglages de Jacobi d'échelon
\note{dans l'isomorphisme qui suit (43), un exposant biffé sur $Q_{0S}^{-1}$.
Dans (44), un $\underline{\omega}_{\underline{J}}$ biffé précède la
parenthèse.}

\page{88}
2 forment les sections d'un torseur sous $(\mathfrak{S}_3)_S$. Donc les
bi-épinglages forment les sections d'un torseur sous le groupe
\struck{constant}
\[
  \text{(45)} \qquad (\mu_4 \times \mathfrak{S}_3)_S
\]
donc on trouve une action \add{à droite} du groupe absolu
$(\mu_4 \times \mathfrak{S}_3)_{S_0}$ sur \struck{\ill{}} le schéma modulaire (37)
$U_{0,3} \times_{S_0} \mu_4^{*}$, qu'il faudrait expliciter, en même temps que
l'action sur
\[
  \text{(46)} \qquad \Pi_0 = \underline{O}_{U_{0,3} \times \mu_4^{*}}(1)
\]
à l'aide de l'action « tautologique » de $\mathfrak{S}_3$ sur
$U_{03} = \Sigma_{\{0,1,\infty\}} \setminus \{0,1,\infty\}_{S_0}$ via
l'action de $\mathfrak{S}_3$ sur $V_{\{0,1,\infty\}}(\underline{O}_{S_0})$.
\[
  \text{(47)} \qquad
  \begin{cases}
    \mathfrak{S}_3 \text{ opère sur } U_{03} \times \mu_4^{*} \text{ via }
    \begin{cases}
      \text{action tautol.\ sur } U_{03} \\
      \text{la \ldots\ } \mathrm{sg}(g) \text{ sur } \mu_4^{*}
    \end{cases} \\[10pt]
    \mu_4 \text{ opère sur } U_{03} \times \mu_4^{*} \text{ via }
    \begin{cases}
      \text{action triviale sur } U_{03} \\
      \text{l'h.\ can.\ } \mu_4 \to \mu_2 = \{\pm 1\}_{S_0} \text{ sur } \mu_4^{*}
    \end{cases}
  \end{cases}
\]
donc $\mathfrak{S}_3 \times \mu_4$ opère sur le fibré $U_{03}$ via l'
\note{devant $U_{03} = \Sigma_{\ldots}$, un $\mathbb{P}$ biffé ; dans
$V_{\{0,1,\infty\}}$, un indice biffé. Le « l'h.\ can. » se lit sans doute
« l'hom.\ can. ». Le mot qui précède $\mathrm{sg}(g)$ commence par
\emph{co-} et n'est pas lu ; nous le rendons « \ldots » dans la formule.}

\page{89}\note{p.~9 de l'auteur.}
action tautologique de $\mathfrak{S}_3$, et sur $\mu_4^{*}$ via la
\ill{} mixte $\mathrm{sg}(g) \cdot \lambda^2$
($g \in \mathfrak{S}_3$, $\lambda \in \mu_4$). Donc le sous-groupe de
$\mathfrak{S}_3 \times \mu_4$ qui opère trivialement sur $\mu_4^{*}$ est
\[
  \text{(48)} \qquad \mathbb{D}'_{3S_0} = (\mathfrak{S}_3)_{S_0}
  \times_{\mu_{2S_0}} (\mu_4)_{S_0}
\]
Si on prend comme base absolue non $S_0$, mais
$\widetilde{S}_0 = (\mu_4^{*})_{S_0}$, de sorte que l'on trouve, sur ces
schémas \add{$/\widetilde{S}_0$} \uncertain{envisagés}
\[
  \begin{cases}
    \mu_{4S}^{*} \text{ muni d'une section canonique } i_S \\
    \mu_{4S} \simeq (\mathbb{Z}/4\mathbb{Z})_S
  \end{cases}
\]
alors le groupe (48) s'identifie au groupe constant sur $S$
\struck{de \ill{}} \add{défini par le groupe ordinaire}
\[
  \text{(49)} \qquad \mathbb{D}'_3 = \mathfrak{S}_3 \times_{\mathbb{Z}/2\mathbb{Z}} (\mathbb{Z}/4\mathbb{Z})
\]
qui est une extension de $\mathbb{D}_3 \simeq \mathfrak{S}_3$ par
$\mathbb{Z}/2\mathbb{Z} \simeq \{\pm 1\}$
\[
  \text{(50)} \qquad 1 \to \{\pm 1\} \to \mathbb{D}'_3 \to \mathbb{D}_3 \to 1 .
\]
\struck{\ill{} action \uncertain{envisagée} de $\mathfrak{S}_3 \times \mu_4$}
Dans le cadre \uncertain{envisagé} \ill{} \uncertain{i.e.} \ill{} travaillant sur
$\widetilde{S}$ \uncertain{i.e.} \ill{} un choix fixé de $i = \sqrt{-1}$, on peut
distinguer
\note{dans (50), le premier terme surcharge un $\mathbb{Z}/2\mathbb{Z}$.
Le début de l'avant-dernière phrase (« \ill{} action \ldots\ de
$\mathfrak{S}_3 \times \mu_4$ ») est biffé et remplacé par l'interligne
« Dans le cadre \ldots ». Nous rendons par $\mathbb{D}$ la lettre D à double
barre de l'auteur.}

\page{90}
les \struck{\ill{}} épinglages de Legendre--Jacobi « directs », i.e.\ tels
que la section de $\mu_4^{*}$ qu'ils définissent soit justement $i$.
\struck{\ill{}} Les \add{champs des} courbes elliptiques relatives
\add{(sur des $S/\widetilde{S}_0$)} munies d'épinglages de Legendre--Jacobi
directs est représentable par $U_{03\widetilde{S}_0}$, sur lequel on fait
opérer le groupe modulaire $\mathbb{D}'_3$, qui est \uncertain{aussi} le
\add{\uncertain{justement}} sous-groupe de
$\mathbb{D}_3 \times \mu_{4\widetilde{S}_0}$ qui fait passer d'un épinglage
direct à un épinglage direct.

Revenant au cas de la base absolue $S_0$, on doit encore expliciter les
opérations de $\mu_4 \times \mathfrak{S}_3$ sur
$\underline{O}_{U_{0,3} \times \mu_4^{*}}(1)$, via l'action tautologique
résultant de l'action de $\mathfrak{S}_3$ sur
$V_{\{0,1,\infty\}}(\underline{O}_{S_0})$ et \add{de $\mathfrak{S}_3$ et}
de $\mu_4$ sur le facteur $\mu_4^{*}$ d'autre part. On désigne l'action
tautologique d'un $(g, \lambda)$ par
$\zeta \in \Gamma\underline{O}_{U_{03}}(1)(A)$
\note{dans $\underline{O}_{U_{0,3} \times \mu_4^{*}}(1)$, un symbole biffé
entre $\underline{O}$ et $(1)$ ; dans la dernière formule, l'indice
$U_{03}$ surcharge un autre. La phrase se poursuit page~91.}

\page{91}\note{p.~10 de l'auteur.}
\add{(où $t : S \to U_{03}$)}, par la \uncertain{substitution} ordinaire
$(g, \lambda) \cdot \zeta$, et on trouve que l'on a, pour l'action
modulaire $T_{g,\lambda}$ sur $\zeta$ :
\[
  \begin{aligned}
    &\text{(51)} \qquad T_{g,\lambda}\, \zeta = \lambda\, (g,\lambda) \cdot \zeta
      && \zeta \in \Gamma\underline{O}_{U_{03}}(1)(t) \\
    &\text{(52)} \qquad T_{g,\lambda}\, \eta = \lambda^2\, (g,\lambda) \cdot \eta
      && \eta \in \Gamma\underline{O}_{U_{0,3}}(2)(t) \\
    &\text{(53)} \qquad T_{g,\lambda}(\tau) = \lambda^2\, \mathrm{sg}(g)\, (g,\lambda) \cdot \tau
      && \tau \in \Gamma\,\underline{O}_{U_{03}}(2)(t)
  \end{aligned}
\]
Quand on se restreint au sous-groupe $\mathbb{D}'_{3S_0}$, pour lequel on a
$\lambda^2 = \mathrm{sg}(g)$, on trouve donc par (51), (52), (53)
\[
  \begin{aligned}
    &\text{(51 bis)} \qquad T_g\, \zeta = \chi_4(g)\, (g \cdot \zeta) \\
    &\text{(52 bis)} \qquad T_g\, \eta = \mathrm{sg}(g^{0})\, (g \cdot \eta) \\
    &\text{(53 bis)} \qquad T_g(\tau) = g \cdot \tau
  \end{aligned}
\]
où $g \in \Gamma(\mathbb{D}'_{3S})$ pour $S$ sur $S_0$,
\[
  \text{(54)} \qquad \chi_4 : \mathbb{D}'_{3S_0} \to \mu_4
\]
le caractère canonique (de sorte que
$\mathbb{D}'_{3\,\mathrm{ab}} \simeq \mu_4$).
\note{la parenthèse initiale « (où $t : S \to U_{03}$) » est écrite au-dessus
de la première ligne et reliée à $\zeta$ par un trait. Dans (51), un mot biffé
suit $\lambda$ ; dans (52), le premier $\eta$ surcharge un $\zeta$ ; dans
(53), le fibré dont $\tau$ est section est surchargé et récrit en dessous,
d'une lecture incertaine (on lit $\underline{O}_{U_{03}}(2)$ et un exposant
biffé). Dans (51 bis), un $\lambda$ biffé suit $T_g$ et un
$\mathrm{sg}(g)$ biffé précède $\chi_4(g)$ ; dans (53 bis), un symbole biffé
précède $g \cdot \tau$. L'exposant $^{0}$ de (52 bis) est lu tel qu'il est
écrit.}

\page{92}
Ces dernières formules sont surtout intéressantes dans le contexte où on
travaille \struck{\ill{}} sur \struck{les} la base absolue $\widetilde{S}_0$ et
avec les structures de Legendre\add{--Jacobi} directes, et où le schéma modulaire
$U_{03} \times_{S_0} \mu_4^{*}$ sur $S_0$ est remplacé par
$U_{03\widetilde{S}_0}$ ($= U_{03} \times_{S_0} \mu_4^{*}$ !) sur
$\widetilde{S}_0$, et le groupe modulaire
$\mathbb{D}_{3S_0} \times_{S_0} \mu_{4S_0}$ est remplacé par le
\add{sous-}groupe $\mathbb{D}'_{3\widetilde{S}_0}$ de
\struck{\ill{}} son image inverse sur $\widetilde{S}_0$.

\note{ici un trait horizontal de l'auteur sépare deux développements.}

Il faut maintenant définir les épinglages de Jacobi d'échelon~4 relatifs
à \uncertain{ceux} d'échelon~2, \struck{i.e.\ des sections} \add{qui sont des
sections \uncertain{(variables)}} de ${}_4E \to {}_2E$ au-dessus de
$\underline{J}_E = ({}_2E)^{*}$, satisfaisant la condition exprimée
\struck{\ill{}} sur chaque fibre géométrique par
$\xi_i + \xi_j + \xi_k = 0$, et montrer comment un tel épinglage définit
canoniquement un épinglage de Legendre.
\note{Dans le passage « relatifs à ceux d'échelon~2 », le mot que nous lisons
« ceux » est très abrégé. Les indices de la relation
$\xi_i + \xi_j + \xi_k = 0$ sont d'une lecture incertaine.}

\page{93}\note{p.~11 de l'auteur.}
Mais avant, un commentaire. Quand on donne $\underline{J}$, rev.\ cubique
étale de $S$, et une section $t$ de $\Sigma_{\underline{J}}^{*}$, dans la
donnée d'une courbe elliptique relative sur $S$, correspondant ainsi aux
données $\underline{J}, t$, équivaut à celle d'un module inv.\ $T$ sur $S$,
et d'un isom (5)
\[
  \text{(*)} \qquad T^{\otimes 2} \simeq T_{\Sigma,t}
\]
Quand on a disposé déjà d'un tel $T$, soit $T_0$, dans la catégorie des $T$
\add{sur $S$} munis de (*) équivaut à celle des $\mu_{2S}$-torseurs
$\underline{\varepsilon}$ sur $S$, moyennant l'équivalence
\[
  \text{(55)} \qquad \underline{\varepsilon} \longmapsto T = T_0 \wedge^{\mu_{2S}} \underline{\varepsilon}
\]
Ici, on a envie de prendre
\[
  \text{(56)} \qquad T_0 = \underline{O}_{\Sigma}(1)(t)^{\natural}
  = \underline{O}_{\Sigma}(1)(t) \wedge^{\mu_4} Q_{0S}
\]
mais on a un iso.\ can.
\[
  T_0^{\otimes 2} \simeq \underline{O}_{\Sigma}(2)(t) \wedge^{\mu_2} \mu_4^{*}
  \simeq T_{\Sigma,t} \wedge^{\mu_2} \bigl(\mu_4^{*} \wedge \underline{\omega}_{\underline{J}}\bigr)
\]
et pour avoir un iso.\ « raisonnable » $T_0^{\otimes 2} \simeq T_{\Sigma,t}$,
il faut se donner un \struck{section} iso
\[
  \alpha : \mu_4^{*} \simeq \underline{\omega}_{\underline{J}}
  \quad \text{i.e.\ une section de } \mu_4^{*} \wedge \underline{\omega}_{\underline{J}}
\]
(ce qui est une des données \uncertain{indiquées} par un épinglage de Legendre
d'une courbe elliptique $E/S$ correspondant à $\underline{J}$\ldots). Une fois
$\alpha$ fixé, on
\note{le premier « dans » de la deuxième phrase est tel qu'écrit. La
référence « (5) » renvoie à une formule d'un lot antérieur ; le repère
« (*) » est un signe de l'auteur. Devant $\alpha$, dans la dernière formule,
un signe biffé.}

\page{94}
trouve l'isom
\[
  k_\alpha : T_0^{\otimes 2} \simeq T_{\Sigma,t}
\]
\add{corr.\ à $\underline{J}, t$} une courbe elliptique $E_\alpha$, munie d'un
\struck{\ill{}} \add{épinglage} de Legendre standard \uncertain{associé} à
$\alpha$ (puisque $T_{E_\alpha,0} \simeq T_0$\ldots). Ceci posé, pour une courbe
elliptique quelconque $E$ sur $S$, corr.\ à $\underline{J}, t$, les
épinglages de Legendre sur $E$ subordonnés à $\alpha$ forment (on l'a vu) les
sections d'un $\mu_2$-torseur $\mathcal{L}_\alpha(E)$ sur $S$, et on voit de
suite qu'on a un isom de $\mu_2$-torseurs
\[
  \text{(57)} \qquad
  \underbrace{\underline{\mathrm{Isom}}_{S,\underline{J}}(E_\alpha, E)}_{%
    \substack{\text{l'indice } \underline{J} \text{ signifie qu'on prend les iso.}\\
    \text{compatibles aux } \underline{J} \simeq {}_2E \simeq {}_2E_\alpha}}
  \xrightarrow{\;\sim\;} \mathcal{L}_\alpha(E),
  \qquad u \longmapsto u(k_\alpha)
\]
Si, au lieu d'une section $\alpha$ de
$\mu_4^{*} \wedge \underline{\omega}_{\underline{J}}$, on disposait d'une
section $\omega$ de $\underline{\omega}_{\underline{J}}$ lui-même, on pourrait
encore définir une courbe elliptique $E_\omega$ canonique corr.\ à
$\underline{J}, t$, et munie d'un « ép.\ de Legendre \uncertain{naïf} »
\add{\uncertain{i.e.} d'un iso.\ $T_{E,0} \simeq \underline{O}_{\Sigma}(1)(t)$
tel que\ldots}), et
\marginal{i.e.\ la courbe elliptique \underline{universelle}, corr.\ à
$\underline{J}, t$ et munie d'un épinglage de Legendre
subordonné à $\alpha$\ldots}
\note{la note marginale est écrite en oblique dans le coin supérieur gauche
et reliée par un trait à « une courbe elliptique $E_\alpha$ » ; sa lecture
est en partie incertaine. Au-dessus de « épinglage », un mot biffé. Avant
(57), une première écriture « $\mathcal{L}_\alpha(E) \leftarrow$ » est biffée.
Le « $\alpha$ » ajouté au-dessus de « section » de l'avant-dernière phrase
est de l'auteur.}

\page{95}
\uncertain{trouvant} une \uncertain{sol.} universelle pour \struck{\ill{}}
\add{de telles} données ($\underline{J}$ torseur sous
$(\mathbb{Z}/3\mathbb{Z})_S$, $t \in \Gamma(\Sigma_{\underline{J}}^{*}/S)$,
$\omega \in \Gamma(S, \underline{\omega}_{\underline{J}})$ correspondant à la
structure de torseur sur $\underline{J}$, et donnant lieu à une
trivialisation de $\underline{\omega}_{\underline{J}}$ sur $S$\ldots). Les
deux points de vue sont équivalents, quand on se fixe une section $i$ de
$\mu_4^{*}$, i.e.\ quand on travaille au-dessus de la base
$\widetilde{S}_0 = \mu_{4S_0}^{*}$ ($S_0 = \operatorname{Spec} \mathbb{Z}[\tfrac{1}{2}]$).

\note{ici un trait horizontal de l'auteur sépare deux développements.}

Explicitons entièrement la notion d'épinglage de Legendre, du point de vue
calculatoire, sur une base réduite à un corps alg.\ clos $k$ pour simplifier,
en décrivant la courbe elliptique $E$ sur $k$ en termes
\[
  \begin{aligned}
    &\text{de } \underline{J} \ (= {}_2E(k)^{*}) \text{ (ens.\ à trois éléments)}, \text{ de la}\\
    &\text{(58)} \quad k\text{-vectoriel } V_{\underline{J}} = \operatorname{Ker}\bigl(k^{\underline{J}} \xrightarrow{\text{somme}} k\bigr) \text{ (d'où} \\
    &\text{(59)} \quad \Sigma_{\underline{J}} = \mathbb{P}^1(V_{\underline{J}}), \text{ avec des } \xi_i \in \Sigma_{\underline{J}}(k) \ (i \in \underline{J}) \text{ distincts,}\\
    &\qquad \Sigma_{\underline{J}}^{*} = \Sigma_{\underline{J}} - \{\xi_i\}_{i \in \underline{J}}\text{)}, \text{ et un } t \in \Sigma_{\underline{J}}^{*} \text{ i.e.\ une droite}\\
    &\text{(60)} \quad F_t \subset V_{\underline{J}},
  \end{aligned}
\]
enfin un $k$-module inv.\ $\underline{t}_E$ \add{($= \underline{t}_E$)} et un
isom
\[
  \text{(61)} \qquad \nu \ (= \nu_E) : \underline{t}_E^{\otimes 2} \simeq T_{\Sigma_{\underline{J}},t} .
\]
Il faut d'abord expliciter calculatoirement $T_{\Sigma_{\underline{J}},t}$
\marginal{$\xi_i = k \cdot (e_j - e_k)$ ou $(e_i)$ est la base canonique de
$V_{\underline{J}}$}
\note{les données « de $\underline{J}$\ldots » et « $k$-vectoriel
$V_{\underline{J}}$\ldots » sont encadrées sur la page ; nous les
disposons en tableau. Dans (58) un symbole biffé suit le signe $=$ ; dans (59)
l'exposant de $\mathbb{P}^1$ surcharge un indice. Dans (61) l'indice $E$ de
$\nu_E$ et celui de $\underline{t}_E$ sont surchargés. La note marginale, écrite
de biais à gauche de (59), est d'une lecture incertaine ; « ou » y est tel
qu'écrit, pour « où ». L'ajout « ($= \underline{t}_E$) » est lu tel quel.}

\page{96}
comme
\[
  \text{(62)} \qquad
  \begin{cases}
    T_{\Sigma_{\underline{J}},t} \simeq \underline{O}_t(2)(\underline{\omega}) \\
    \text{avec } \underline{O}_t(2) \overset{\mathrm{def}}{\simeq} \underline{O}_t(1)^{\otimes 2},
    \quad \underline{O}_t(1) \simeq V_{\underline{J}}/F_t \simeq \check{F}_t(\underline{\omega})
  \end{cases}
\]
de sorte que $\nu_E$ s'interprète comme un isom
\[
  \text{(63)} \qquad \nu : \underline{t}_E^{\otimes 2} \simeq F_t^{\otimes(-2)}(\underline{\omega})
  \qquad \text{ou aussi comme} \qquad
  \check{\nu} : \underline{t}^{\otimes -2} \simeq F_t^{\otimes 2}(\underline{\omega})
\]
(NB $\underline{\omega} = \underline{\omega}_{\underline{J}}$, un des deux ordres
circulaires sur $\underline{J}$). Les données de base sont donc,
pour nous, $\bigl[\underline{J}$ ($\mapsto \underline{\omega} = \underline{\omega}_{\underline{J}}$,
\ill{} et $V_{\underline{J}} = V_{\underline{J}}(k)$),
$F_t \hookrightarrow V_{\underline{J}}$, $\underline{t}$,
$\nu : \underline{t}^{\otimes 2} \simeq F_t^{\otimes(-2)}(\underline{\omega})\bigr]$.

Ceci posé, la donnée d'un épinglage de Legendre s'explicite comme un isom
\[
  \text{(64)} \qquad k : \underline{t} \xrightarrow{\;\sim\;} \underline{O}_t(1)^{\natural}
\]
i.e.\ comme une famille d'isomorphismes
\[
  \text{(65)} \qquad
  \begin{cases}
    (k_i)_{i \in \mu_4^{*}}, \quad k_i : \underline{t} \xrightarrow{\;\sim\;} \underline{O}_t(1) \simeq \check{F}_t(\underline{\omega}) \\
    k_{i'} = i\, k_i \qquad \text{si } i' = 1/i \text{ est l'autre élément de } \mu_4^{*}
  \end{cases}
\]
ou encore, dans familles
\[
  \text{(66)} \qquad
  \begin{cases}
    (\check{k}_i)_{i \in \mu_4^{*}}, \quad \check{k}_i : \check{\underline{t}} \simeq F_t(\underline{\omega}) \\
    \check{k}_{i'} = \tfrac{1}{i}\, \check{k}_i ,
  \end{cases}
\]
satisfaisant à une certaine condition :
\[
  \text{(67)} \qquad
  \begin{cases}
    k_i^{\otimes 2} : \underline{t}^{\otimes 2} \xrightarrow{\;\sim\;} \check{F}_t^{\otimes(-2)} \text{ est de la forme} \\
    k_i^{\otimes 2}(\delta^{\otimes 2}) = \nu(\delta^{\otimes 2}) \wedge \omega_i
  \end{cases}
\]
où $\omega_i \in \underline{\omega}$ est un élément convenable de
$\underline{\omega}$, bien déterminé non par $k, i$, en fait par $k_i$.
\marginal{$\mu_4^{*} = \mu_4^{*}(k)$, racines d'ordre $4$ de $k^{*}$}
\note{dans (63), l'accent de $\check{F}_t$ et l'exposant $\otimes(-2)$ sont
surchargés. Le mot illisible de la liste des données de base est un
symbole barré suivi de « et ». Dans (67), l'accent sur $F_t$ est un
circonflexe surchargé ; la lettre que nous lisons $\delta$ (comme page~97,
où $\delta \in \underline{t}^{*}$) est tracée ici de façon peu nette, et n'est pas
introduite sur la page.
L'indice $i$ de $k_i$, $\omega_i$ est écrit avec un point et désigne un
élément de $\mu_4^{*}$.}

\page{97}
On aura donc
\[
  \text{(68)} \qquad
  \begin{cases}
    \omega_{i'} = -\omega_i \\
    \text{i.e.\ } i \mapsto \omega_i \text{ est une bijection} \\
    \qquad \alpha = \alpha_k : \mu_4^{*} \xrightarrow{\;\sim\;} \underline{\omega}
  \end{cases}
\]
On va expliciter la donnée des $k$, \add{donnant un préalable}
\struck{$\alpha : \mu_4^{*} \simeq \underline{\omega}$} associé à un isom.\
$\alpha = (i \mapsto \omega_i) : \mu_4^{*} \simeq \underline{\omega}$,
via les $\check{k}_i$. On aura, pour $\delta \in \underline{t}^{*}$
\[
  \text{(69)} \qquad \check{k}_i(\check{\delta}) = \underbrace{\psi_i(\delta)}_{\in F_t} \wedge \omega_i
  = \underbrace{\psi^{\omega_i}(\delta)} \wedge \omega_i ,
\]
où
\[
  \text{(70)} \qquad \bigl(\psi^{\omega}(\delta)\bigr)_{\omega \in \underline{\omega}} \in (F_t)^{\underline{\omega}}
\]
est une famille \struck{\ill{}} de deux éléments de $F_t$, indexée par
$\underline{\omega}$. La condition \add{(65) ou} (66) s'explicite comme
\[
  \underbrace{\psi^{\omega_{i'}}(\delta) \wedge \omega_{i'}}_{\substack{\psi^{\omega'}(\delta) \wedge \omega' \\ = -\psi^{\omega'}(\delta) \wedge \omega}}
  = \frac{1}{i}\, \underbrace{\psi^{\omega_i}(\delta) \wedge \omega_i}_{\psi^{\omega}(\delta) \wedge \omega}
  \qquad \text{i.e.\ } \psi^{\omega'}(\delta) = -\frac{1}{i}\, \psi^{\omega}(\delta)
\]
\struck{\ill{}} si \add{on pose} $\omega \overset{\mathrm{def}}{=} \omega_i$,
$\omega' = $ autre él.\ de $\underline{\omega}$, i.e.
\[
  \text{(71)} \qquad \psi^{\omega'}(\delta) = i_\omega\, \psi^{\omega}(\delta)
  \qquad \text{où } i_\omega \overset{\mathrm{def}}{=} \alpha^{-1}(\omega)
\]
On doit avoir de plus
\[
  \text{(72)} \qquad \psi^{\omega}(\lambda\delta) = \frac{1}{\lambda}\, \psi^{\omega}(\delta) ,
\]
\struck{qui ex} et enfin il faut exprimer (67),
\marginal{\underline{NB} Comme la donnée d'une \ldots\ dans l'iso.\
$\alpha : \mu_4^{*} \simeq \underline{\omega}$, les deux relations (67)
\ill{} \ill{}\ldots}
\note{sous « $-\frac{1}{i}$ » de la condition, l'auteur écrit par une
accolade « $i$ » : $-1/i = i$ dans $\mu_4^{*}$, d'où (71). Dans (68), un
mot biffé en tête, et un autre après « bijection » ; l'indice de $\alpha_k$ surcharge une autre lettre. La
note marginale, écrite de biais en haut à gauche, n'est lue qu'en partie.
Le $\check{\delta}$ de (69) porte un accent sur la lettre et un point devant
$\check{k}$.}

\page{98}
qui précisément sous forme contragrédiente
\[
  \check{k}_i^{\otimes 2}(\check{\delta}^{\otimes 2}) = \check{\nu}(\check{\delta}^{\otimes 2}) \wedge \omega_i ,
\]
en termes des $\psi^{\omega}(\delta)$, comme
\[
  \text{(73)} \qquad \psi^{\omega}(\delta)^{\otimes 2} = \check{\nu}(\check{\delta}^{\otimes 2}) \wedge \omega .
\]
On trouve ainsi le

\underline{Lemme} Les épinglages de Legendre de données
$\bigl(\underline{J}, F_t \subset V_{\underline{J}}(k), \underline{t},
\check{\nu} : \check{\underline{t}}^{\otimes 2} \simeq F_t^{\otimes 2}(\underline{\omega})\bigr)$,
s'identifiant aux systèmes
$\bigl(\alpha, (\psi^{\omega}(\delta))_{\omega \in \underline{\omega}}\bigr)$, où
$\alpha : \mu_4^{*} \simeq \underline{\omega}$, et où
$(\psi^{\omega}(\delta))_{\omega \in \underline{\omega},\, \delta \in \underline{t}^{*}}
\in (F_t)^{\underline{\omega} \times \underline{t}^{*}}$
\struck{satisfaisant} \add{\uncertain{produit} $\underline{\omega} \times \underline{t}^{*}$}
est une famille indexée par l'\uncertain{ens.}\ \ill{} (notons
$\underline{\omega} = \underline{\omega}_{\underline{J}}$), satisfaisant les
conditions (71) (72) (73).

\struck{\ill{}} Nous allons maintenant expliciter la donnée d'une famille
\add{à deux indices}
$(\psi^{\omega}(\delta))_{\omega \in \underline{\omega},\, \delta \in \underline{t}^{*}}$,
à val.\ dans $F_t$, en termes des \ill{}. Moyennant
\[
  F_t \subset V_{\underline{J}} = \operatorname{Ker}\bigl(k^{\underline{J}} \xrightarrow[\text{somme}]{} k\bigr) ,
\]
on voit donc qu'elles s'expriment comme une famille à trois indices
d'éléments de $k$ :
\[
  \text{(74)} \qquad \bigl(\psi_i^{\omega}(\delta)\bigr)_{\substack{i \in \underline{J} \\ \omega \in \underline{\omega} \\ \delta \in \underline{t}^{*}}},
  \qquad \psi_i^{\omega}(\delta) \in k
\]
Le fait que la famille d'éléments de $k^{\underline{J}}$
\marginal{\underline{NB} Moyennant (71) (\uncertain{ou} (72)), les deux
relations (73), pour les deux él.\ $\omega \in \underline{\omega}$, sont
équivalentes, il suffit \ill{}\ldots}
\note{« Lemme » est souligné. Dans l'énoncé, l'ajout interligne, sous un
« satisfaisant » biffé, se lit « pr.\ \ill{} $\underline{\omega} \times
\underline{t}^{*}$ » et est lui-même souligné d'un trait qui le barre peut-être.
La note marginale, serrée en haut à droite, est surchargée dans ses deux
dernières lignes. Le mot qui suit « en termes des » n'est pas lu ;
« Moyennant » est une lecture du mot qui ouvre la ligne suivante.}

\page{99}
indexée par $\underline{\omega} \times \underline{t}^{*}$ soit une famille
d'éléments de $V_{\underline{J}}(k)$ s'exprime par la formule
\[
  \text{(75)} \qquad \sum_{i \in \underline{J}} \psi_i^{\omega}(\delta) = 0
\]
les \struck{\ill{}} conditions (71) et (72) s'expriment ensuite en termes des
$\psi_i^{\omega}(\delta)$, comme
\[
  \begin{aligned}
    &\text{(76)} \qquad \psi_i^{\omega'}(\delta) = i_\omega\, \psi_i^{\omega}(\delta) \\
    &\text{(77)} \qquad \psi_i^{\omega}(\lambda\delta) = \lambda^{-1}\, \psi_i^{\omega}(\delta) .
  \end{aligned}
\]
Il reste à exprimer le fait que les éléments $\psi^{\omega}(\delta) \in V_{\underline{J}}(k)$
définis par les $(\omega, \delta) \in \underline{\omega} \times \underline{t}^{*}$ sont
dans $F_t$ --- moyennant (76), (77) il suffit de l'exprimer pour un seul
couple $(\omega, \delta)$ ---, et la condition (73).
\struck{Si on définit $F_t$ \ill{} comme \ill{} deux \ill{} \ill{} d'une forme
linéaire sur $V_{\underline{J}}$ \ill{}, on écrira}
\struck{(77) \quad $\sum_{i \in \underline{J}} \psi_i^{\omega}(\delta)\, e_i$ ---}
\[
  F_t = \Bigl\{ \sum \lambda_i e_i \Bigm| \sum \lambda_i = 0,\
  \sum c_i \lambda_i = 0 \Bigr\}
\]
[\underline{NB} L'élément $(c_i)_{i \in \underline{J}}$ de $k^{\underline{J}}$ est
déterminé modulo multiplication par $k^{*}$ et addition d'un élément
diagonal], \struck{cette} la condition $\psi^{\omega}(\delta) \in F_t$ s'exprime
par la condition calculatoire
\[
  \text{(78)} \qquad \sum_{i \in \underline{J}} c_i\, \psi_i^{\omega}(\delta) = 0 .
\]
Il reste \add{ensuite} à exprimer (73)
\marginal{$(e_i)_{i \in \underline{J}}$ base canonique de $k^{\underline{J}}$}
\note{le passage biffé entre « la condition (73) » et la formule donnant
$F_t$ occupe deux lignes et une formule numérotée (77), elle aussi biffée ; le
numéro (77) est donc employé deux fois sur la page. Devant l'accolade de
$F_t$, et devant $k^{\underline{J}}$ dans le NB, un symbole biffé ; dans la
formule de $F_t$, un facteur biffé après $\sum c_i \lambda_i$.}


\section{IV Descriptions explicites par fonctions $\wp$ de Weierstrass}
\note{titre écrit de la main de l'auteur en haut d'une feuille (page~100 de
l'archive) par ailleurs blanche ; il annonce la suite, qui commence au lot
suivant. Le « $\wp$ » est surchargé sur le mot « fonctions ».}

\end{document}
